\documentclass[11pt]{amsart}
\usepackage{amsmath,amssymb,amsthm}
\usepackage[margin=1.05in]{geometry}
\usepackage{booktabs}
\usepackage{enumitem}
\usepackage{hyperref}
\emergencystretch=2em

\newtheorem{theorem}{Theorem}[section]
\newtheorem{lemma}[theorem]{Lemma}
\newtheorem{proposition}[theorem]{Proposition}
\newtheorem{corollary}[theorem]{Corollary}
\theoremstyle{definition}
\newtheorem{definition}[theorem]{Definition}
\newtheorem{remark}[theorem]{Remark}
\newtheorem{question}[theorem]{Question}

\newcommand{\Ev}{\mathrm{Ev}}
\newcommand{\Od}{\mathrm{Od}}
\newcommand{\tr}{\rightarrowtail}   % transfer of one end
\newcommand{\bl}{\mathrm{bl}}

\title[Hub transfers at deletion degree six]{Hub transfers at deletion degree six:\\
closures and obstructions for the clique core with full exterior vertices}
\author{Working draft for the Gallai path-decomposition programme}
\date{26 November 2026}

\begin{document}

\begin{abstract}
The open case of the deletion-degree-six rooted test (T(6)-i) asks, for a connected graph $B$ of
order $2k$ with exactly six even vertices $U$ and a prescribed $r\in U$, for a partition of $E(B)$
into at most $k$ paths with two ends at $r$. When $B[U]=K_6$, by six-clique Thm.~6~\cite{6CL} the only unresolved
instances have an exterior vertex adjacent to all of $U$ (a \emph{full} vertex). We do not close this case. We prove
four kinds of results, each by an explicit short sequence of Fan addible-edge transfers starting
from a Lov\'asz seed, so that each is directly expressible with the formally verified lemmas
\texttt{fan\_single}, \texttt{fan\_grouped} and \texttt{realisable\_evenSet\_erase}.
(1) A \emph{terminal-step classification}: if all held edges lie in $U$ together with the full
vertices, the last Fan step of any such certificate can only end at one of three ``hub'' profiles.
It can never end at the standard profile with three doubled core vertices. This is the precise
form of the obstruction behind the ``off-by-one'' failures of single-edge returns.
(2) \emph{Root-hub closures}: if $r$ is adjacent to all other even vertices, has degree at least
$8$, and has an odd neighbour $y$ of degree at least $3$ with at most two neighbours in $U-r$, then
$r$ passes. Similar small-contact witnesses (a triangle at $r$, a pair of resources,
or a relay path to a vertex with no core neighbour) give closures, under analogous hypotheses, at
the root positions of every residual core shape $D\in\{\emptyset,K_2,P_3,K_3,C_4\}$.
(3) A \emph{full-hub closure}: a full vertex carrying two ``quiet batteries'' (triangles through
it whose parked vertex has no core neighbour) closes the test at every root.
(4) Uniform-in-$d$ versions of (2) and (3) for every deletion degree $d$. They include the first
closures of the open test T(8)-i at roots adjacent to all other even vertices.
We also show by an exhaustive (small) search that no single-stage program inserts a pair of full
vertices in the six-clique stage framework. This explains why that method stops at five contacts.
All certificates were replayed on concrete graphs by an independent implementation of Fan's
construction. The relay lemma, the root hubs (2) of Theorem~\ref{thm:roothub}, the full hub (3) and both
uniform theorems (4) are also formalized in Lean, with no fixed $d$ (Remark~\ref{rem:lean}).
\end{abstract}

\maketitle

\section{Setting and summary}\label{sec:intro}

Graphs are finite and simple. A \emph{path} is a nonempty simple path. A \emph{partition} of a graph
is a partition of its edge set into paths. For a partition $D$ and a vertex $v$, let $b_D(v)$ be the
number of path ends at $v$. Then
\begin{equation}\label{eq:parity}
b_D(v)\equiv d(v)\pmod 2,\qquad \sum_v b_D(v)=2|D|.
\end{equation}
Write $\Od(G)$ and $\Ev(G)$ for the odd-degree and even-degree vertices.

\subsection*{The test} Throughout, $B$ has even order $2k$ and no isolated vertex (seeds need
this, and so does the count $|\Od(B)|=2k-6$), $U=\Ev(B)$ has six vertices, $r\in U$
and $T=U-r$. For a vertex $x$ write $S(x)=N(x)\cap U$. \emph{$r$ passes} if $B$ has a partition into at most $k$ paths with $b(r)\ge 2$. By
the private-leaf equivalence and Marked Reduction Thm.~114, Gallai's conjecture at deletion
degree~6 is equivalent to: every $r$ passes in every $B$ of the apex class. In the apex class,
$J=B+z$ with $N(z)=U$ is $3$-connected, $\delta(B)\ge 6$, and so on; see the status map
\texttt{research/notes/reader-six-even-program.md}. With $D=K_6-B[U]$, the shape theorem leaves
$D\in\{\emptyset,K_2,P_3,K_3,C_4\}$ with restricted roots. Six-clique Thm.~6 closes $D=\emptyset$
unless some exterior vertex is \emph{full}, i.e.\ adjacent to all of $U$.

A partition into $k$ paths with $b=1$ on $\Od(B)$ has $2k-(2k-6)=6$ ends on $U$. If moreover $b\le2$
on $U$, exactly three core vertices are \emph{doubled} ($b=2$) and three are \emph{blank} ($b=0$).
We call this the \emph{standard profile}. A vertex with $b=0$ in the current partition is
\emph{blank}.

\subsection*{Results} Table~\ref{tab:summary} lists what is proved here. ``New'' is relative to the
collection as summarised in the status map and audit A6.

\begin{table}[ht]
\small
\begin{tabular}{p{2.3cm}p{8.6cm}p{2.9cm}}
\toprule
Result & Statement (informal) & Status\\
\midrule
Thm.~\ref{thm:terminal}, \ref{thm:terminalW} & if all held edges are local, the last Fan step ends at a hub profile, never at the standard profile & proved\\
Thm.~\ref{thm:roothub} & (i) $r$ adjacent to all of $T$ (so $D=\emptyset$, or $r\notin V(D)$), with an odd $y\in N(r)$ and $|N(y)\cap T|\le2$; (ii), (iii) $r$ has one or two core non-neighbours and a triangle $rab$ with a small-contact $a$ & proved; new for $|S(y)|=3$ and for (ii), (iii); Lean\\
Thm.~\ref{thm:tworesources} & $D=\emptyset$: any odd $y\in N(r)$ plus a triangle $rab$, under mild contact conditions & proved\\
Thm.~\ref{thm:relay} & $D=\emptyset$: a relay path from a non-$T$-full neighbour of $r$ to a vertex with no core neighbour & proved\\
Thm.~\ref{thm:fullhub} & a full vertex with two quiet batteries closes every root & proved; Lean\\
Thm.~\ref{thm:general} & hub closures for every deletion degree $d$, incl.\ T(8)-i at universal roots & proved; Lean\\
Prop.~\ref{prop:stage} & no single-stage full-pair program in the six-clique framework & exhaustive small search\\
\midrule
Q2 & $D=\emptyset$ with full vertices, in general & \textbf{open}; residual in Cor.~\ref{cor:residual}\\
\bottomrule
\end{tabular}
\caption{Summary.}\label{tab:summary}
\end{table}

\subsection*{Inputs} Only three facts are used.
\begin{enumerate}[label=(\Alph*)]
\item Lov\'asz~\cite{L68}: a graph on $N$ vertices with at most one even vertex has a partition into
$\lfloor N/2\rfloor$ paths. In Lean this is the theorem \texttt{lovaszStatement} (proved from Fan's lemma,
\texttt{manuscripts/lovasz-fan}), so the formalized results of this paper are unconditional.
\item Endpoint-structure Lemma~3~\cite{ES}, a consequence of (A). If $G$ has even order and no
isolated vertex and $u\in\Ev(G)$, then $G$ has a partition with $b=2$ on $\Ev(G)-u$, $b=0$ at $u$ and
$b=1$ on $\Od(G)$. In Lean this is \texttt{realisable\_evenSet\_erase'}.
\item Fan's Lemmas 3.3 and 3.4~\cite{Fan}, stated below. In Lean these are
\texttt{fan\_grouped} and \texttt{fan\_single}.
\end{enumerate}

\section{Transfers, seeds and relays}\label{sec:tools}

\begin{lemma}[Fan~\cite{Fan}]\label{lem:fan}
Let $D$ be a partition of a graph $H$, let $v\in V(H)$, and let $S$ be a set of vertices that are
not adjacent to $v$ in $H$, with $v\notin S$. Write
$\bl_H(v)=|\{w\in N_H(v):b_D(w)=0\}|$. If $\sum_{s\in S}b_D(s)>\bl_H(v)$, then for some $s\in S$
the graph $H+sv$ has a partition $D'$ with $|D'|=|D|$, $b_{D'}(v)=b_D(v)+1$,
$b_{D'}(s)=b_D(s)-1$ and $b_{D'}=b_D$ elsewhere. If $S=\{s\}$ the conclusion holds for this $s$.
\end{lemma}

We write $s\tr v$ (``$s$ gives an end to $v$'') for one such insertion. For $S=\{s\}$ it is
\emph{single}; otherwise it is \emph{grouped}, and the donor is not under our control. In a
certificate every quantity we use is either the current end vector or an upper bound on
$\bl_H(v)$. We record each step as a pair (supply, bound).

\begin{lemma}[Seeds]\label{lem:seeds}
Let $G$ have even order and no isolated vertex.
\begin{enumerate}[label=(\alph*)]
\item If $\Ev(G)=\emptyset$, then $G$ has a partition into $|G|/2$ paths with $b=1$ everywhere.
\item If $\Ev(G)=\{u,w\}$, then $G$ has a partition into $|G|/2$ paths with $b(u)=2$, $b(w)=0$ and
$b=1$ on $\Od(G)$.
\end{enumerate}
\end{lemma}
\begin{proof}
(a) Apply (A). The partition has at most $|G|/2$ paths, and every vertex is odd and so needs at
least one end. Counting ends by \eqref{eq:parity} gives exactly one end at each vertex. (b) is (B)
with $\Ev(G)-w=\{u\}$. The number of paths is $(|G|-2)/2+1$.
\end{proof}

A certificate in this paper always has the same shape. Choose a set $H$ of \emph{held} edges so
that $B-H$ satisfies Lemma~\ref{lem:seeds}. Take the seed. Then re-insert the held edges one at a
time by Lemma~\ref{lem:fan}. The path count never changes, so a certificate that ends with all
edges present and $b(r)\ge2$ shows that $r$ passes.

\begin{lemma}[Relay]\label{lem:relay}
Let $D$ be a partition of $H$ and let $v_0v_1\cdots v_m$ $(m\ge1)$ be a path whose edges are not
in $H$. Put $H_i=H+v_0v_1+\cdots+v_{i-1}v_i$.
\begin{enumerate}[label=(\alph*)]
\item \emph{Forward relay.} Suppose $b(v_i)=1$ for $1\le i<m$ and $b(v_0)>\bl_{H}(v_1)$.
Suppose also that, for $2\le i\le m$, the vertex $v_i$ has at most one blank neighbour in
$H_{i-1}$, counted for the partition current at that moment. Then $H_m$ has a partition with the
same number of paths, $b(v_0)$ decreased by one, $b(v_m)$ increased by one, and all other counts
unchanged.
\item \emph{Carry.} Let $B_0$ be the set of blank vertices of $D$. Suppose $b(v_i)=1$ for
$1\le i\le m$ and that none of $v_0,\dots,v_{m-1}$ is adjacent in $H_m$ to a vertex of $B_0$.
Then the transfers $v_1\tr v_0,\ v_2\tr v_1,\dots,v_m\tr v_{m-1}$ succeed. Afterwards $b(v_0)$ is
increased by one, $b(v_m)=0$, and all other counts are unchanged.
\end{enumerate}
\end{lemma}
\begin{proof}
(a) Apply single transfers $v_0\tr v_1,\ v_1\tr v_2,\dots$ in turn. The first needs
$b(v_0)>\bl_H(v_1)$. After $v_{i-2}\tr v_{i-1}$ the vertex $v_{i-1}$ has two ends, so
$v_{i-1}\tr v_i$ needs $\bl_{H_{i-1}}(v_i)\le 1$. It then leaves $v_{i-1}$ with one end again.
(b) The first transfer $v_1\tr v_0$ has supply $1$, and $v_0$ has no blank neighbour. After
$v_i\tr v_{i-1}$ the vertex $v_i$ has no end, and the blank vertices are $B_0\cup\{v_i\}$. The
next transfer $v_{i+1}\tr v_i$ has supply $1$, and $v_i$ is adjacent to no vertex of $B_0$.
\end{proof}

A carry followed later by a single transfer $v_0\tr v_m$ along a further edge $v_0v_m$ is a
\emph{battery}. While $v_m$ is \emph{parked} (blank), $v_0$ holds one extra end, and the final
transfer returns it.

\section{The terminal-step classification}\label{sec:terminal}

The next theorem explains why returns of local edges fail ``by one'' in the clique case. Examples
are connections \S J and false-twin-return Prop.~4~\cite{FTR}.

\begin{theorem}[Terminal step]\label{thm:terminal}
Let $B$ have order $2k$, $U=\Ev(B)$, $|U|=6$, $B[U]=K_6$. Let $F\subseteq\Od(B)$ be a set of
vertices each adjacent to all of $U$, and put $L=U\cup F$. Let $e=uv$ be an edge with $u,v\in L$.
Let $D'$ be a partition of $B-e$ into $k$ paths with
$b_{D'}(u)>\bl_{B-e}(v)$. Let $D$ be the partition of $B$ produced by Lemma~\ref{lem:fan}
(single, $u\tr v$). If $b_D(w)\ge2$ for some $w\in U$, then exactly one of the following holds.
\begin{enumerate}[label=($\alpha$)]
\item[($\alpha$)] $u\in U$, $b_D(u)=4$, there is exactly one $w\in U-u$ with $b_D(w)=2$, $b_D=0$ on
the other four vertices of $U$, $b_D=1$ on $\Od(B)$, and $v\in\{w\}\cup F$.
\item[($\beta$)] $u\in U$, $b_D(u)=6$, $b_D=0$ on $U-u$, $b_D=1$ on $\Od(B)$, and $v\in F$.
\item[($\gamma$)] $u\in F$, $b_D(u)=5$, exactly one $w\in U$ has $b_D(w)=2$, $b_D=0$ on $U-w$,
$b_D=1$ on $\Od(B)-u$, and $v\in\{w\}\cup(F-u)$.
\end{enumerate}
In particular $D$ never has the standard profile.
\end{theorem}

\begin{proof}
Write $b=b_D$. Then $b_{D'}=b+e_u-e_v$, so $b(v)\ge1$. Define the excess $\varepsilon(w)=b(w)$ for
$w\in U$ and $\varepsilon(w)=b(w)-1$ for $w$ odd. By \eqref{eq:parity} every $\varepsilon(w)$ is
even and nonnegative, and $\sum_w\varepsilon(w)=2k-(2k-6)=6$. Let $Z=\{w\in U:b(w)=0\}$ and
$\beta=|Z|$. Every vertex of $U-Z$ has $b\ge2$, so $2(6-\beta)\le 6$ and $\beta\ge3$.

The vertex $v$ lies in $L$, so it is adjacent in $B$ to every vertex of $U-v$, and $v\notin Z$. A
vertex $w\in Z-\{u\}$ has $b_{D'}(w)=0$ and is adjacent to $v$ in $B-e$. Hence
$\bl_{B-e}(v)\ge|Z-\{u\}|$, and Fan's hypothesis gives
\begin{equation}\label{eq:term}
b(u)+1=b_{D'}(u)>|Z-\{u\}|.
\end{equation}

\emph{Case $u\in Z$.} Then \eqref{eq:term} reads $1>\beta-1$, so $\beta\le1$, which contradicts
$\beta\ge3$.

\emph{Case $u\in U-Z$.} Now $b(u)\ge\beta$. The other $5-\beta$ vertices of $U-Z$ each carry excess
at least $2$, so $b(u)+2(5-\beta)\le6$, that is $b(u)\le2\beta-4$. Together with $b(u)\ge\beta$
this gives $\beta\ge4$. If $\beta=4$, then $b(u)=4$, the one remaining vertex $w$ of $U-Z-u$ has
$b(w)=2$, and no excess is left. That is ($\alpha$). If $\beta=5$, then $b(u)\ge5$ is even, so
$b(u)=6$ and all excess is at $u$. That is ($\beta$). In both cases $v\ne u$ needs $b(v)\ge1$ with
no excess, so $v=w$ or $v\in F$ in ($\alpha$), and $v\in F$ in ($\beta$).

\emph{Case $u\in F$.} Write $b(u)=1+2a$. Then \eqref{eq:term} gives $2+2a>\beta$. The excess count
gives $2a+2(6-\beta)\le6$, so $a\le\beta-3$. Hence $2+2(\beta-3)>\beta$, that is $\beta\ge5$. The
hypothesis that some $w\in U$ has $b(w)\ge2$ forces $\beta=5$. Then $a\ge2$ from $2+2a>5$, and
$a\le2$, so $b(u)=5$, $b(w)=2$ for the unique $w\in U-Z$, and there is no other excess. The receiver
must then be $w$ or a vertex of $F-u$.
\end{proof}

\begin{theorem}[Terminal step, leaf form]\label{thm:terminalW}
Let $W=B+\ell$, where $\ell$ is a new leaf at $r$, so that $|W|=2k+1$ and $\Ev(W)=T$. With
$B[U]=K_6$ and $F$ as above, put $L=T\cup\{r\}\cup F$. If the last edge $uv$, with $u,v\in L$, is
inserted by a single Fan step into a partition of $W-uv$ into $k$ paths, then either
\begin{enumerate}[label=($\alpha_W$)]
\item[($\alpha_W$)] $u\in T$, $b(u)=4$, $b=0$ on $T-u$, $b=1$ on $\Od(W)$, and $v\in\{r\}\cup F$; or
\item[($\gamma_W$)] $u\in\{r\}\cup F$, $b(u)=5$, $b=0$ on $T$, $b=1$ on $\Od(W)-u$, and
$v\in(\{r\}\cup F)-u$.
\end{enumerate}
\end{theorem}
\begin{proof}
The argument is the same, with total excess $2k-(2k-4)=4$ and $|T|=5$. One gets $\beta\ge3$ and
\eqref{eq:term}. The case $u\in Z$ is impossible. For $u\in T-Z$ one gets $b(u)\le2\beta-4$ and
$b(u)\ge\beta$, hence $\beta=4$ and $b(u)=4$. For $u$ odd one gets $\beta=5$ and $b(u)=5$.
\end{proof}

\begin{corollary}\label{cor:nonlocal}
Consider any certificate (\S\ref{sec:tools}) for the rooted test at $D=\emptyset$ that ends at the
standard profile, or any certificate for the floor bound of $W$. Its last inserted edge has an end
outside $L$, unless the certificate ends at a hub profile ($\alpha$)--($\gamma$), resp.\
($\alpha_W$), ($\gamma_W$). (The last step of a certificate is always a single transfer: every
donor offered in a grouped step must be the far end of a held edge at the receiver, and only one
held edge is left. So Theorems~\ref{thm:terminal} and~\ref{thm:terminalW} apply to it.) In particular, returning a single edge $rx$ at a full $x$ from a
partition of $B-rx$ rooted at $x$ is never justified by Fan's inequality when the outcome is
standard. This is the numerical obstruction of false-twin-return Prop.~4.
\end{corollary}

\begin{remark}[Induction on the number of full vertices]\label{rem:induction}
Deleting the twelve spokes of two full vertices $x,x'$ keeps the class: $U$ stays even and
$B[U]=K_6$, and there are two fewer full vertices. Deleting the single edge $rx$ also keeps the
class, with root $x$. An induction would take a rooted partition of the smaller graph and return
the deleted edges by Fan transfers. Every returned edge is local, so by Theorem~\ref{thm:terminal}
the last return can only produce a hub profile. A standard rooted profile is never reached this
way. Any induction on full vertices therefore needs either a non-Fan final step or non-local held
edges, and the hub theorems below provide the second.
\end{remark}

\begin{remark}[Hub capacity]\label{rem:capacity}
The hub profiles need a vertex with $4$--$6$ ends. Fan receipts cannot produce these from the
clique alone. Let a vertex $h$ have no blank neighbour, and let it receive successively from $m$
non-neighbours, each with exactly one end, where each spent donor becomes a blank neighbour of
$h$. After $j$ receipts the remaining supply is $m-j$ and $h$ has $j$ blank neighbours, so at most
$\lceil m/2\rceil$ receipts are possible. With $m=5$, the hub $r$ (or a full $x$) collects three
ends from $T$. A hub that must then donate to receivers adjacent to all three new blanks needs
supply at least $4$ at each donation. The \emph{root hub} $r$ makes two such donations and ends
with two ends. The \emph{full hub} makes three and ends with one. Counting ends, the root hub is
short by exactly one unit and the full hub by exactly two. The resources of
\S\S\ref{sec:roothub}--\ref{sec:fullhub} supply these units through non-local edges, as
Corollary~\ref{cor:nonlocal} requires.
\end{remark}

\section{Root hubs}\label{sec:roothub}

In this section $Q=T-N(r)$ is the set of core non-neighbours of $r$. Note $|Q|=\deg_D(r)$. The
residual root positions of the shape theorem have $|Q|=0$ (all roots for $D=\emptyset$, and roots
off the edge for $D=K_2$), $|Q|=1$ (the roots on the edge for $D=K_2$), and $|Q|=2$ (the middle of
$P_3$, a vertex of $K_3$, a vertex of $C_4$). The theorems below do not assume a shape. They use
only $|Q|$.

\begin{theorem}[Root hub]\label{thm:roothub}
Let $B$ have even order $2k$, $U=\Ev(B)$, $|U|=6$, $r\in U$, $T=U-r$, $Q=T-N(r)$.
\begin{enumerate}[label=(\roman*)]
\item If $Q=\emptyset$, $d(r)\ge8$, and $r$ has an odd neighbour $y$ with $d(y)\ge3$ and
$|N(y)\cap T|\le2$, then $r$ passes.
\item If $Q=\{q\}$, $q$ is not isolated, $d(r)\ge8$, and $r$ has two adjacent odd neighbours $a,b$ with
$|N(a)\cap T|\le2$ and $bq\notin E(B)$, then $r$ passes.
\item If $Q=\{q_1,q_2\}$, neither $q_i$ is isolated, and $r$ has two adjacent odd neighbours $a,b$ with
$|N(a)\cap N_T(r)|\le1$ and $bq_2\notin E(B)$, then $r$ passes.
\end{enumerate}
In each case the partition has exactly $k$ paths, one end at every odd vertex, and three doubled
vertices of $U$, one of which is $r$.
\end{theorem}

\begin{proof}
(i) Hold $H=\{rt:t\in T\}\cup\{ry\}$. In $B_0=B-H$ the vertex $r$ has even degree $d(r)-6\ge2$, $y$
has even degree $d(y)-1\ge2$, the vertices of $T$ are odd, and nothing else changed. So
$\Ev(B_0)=\{r,y\}$ and $B_0$ has no isolated vertex. Lemma~\ref{lem:seeds}(b) gives a seed with
$k$ paths, $b(r)=2$, $b(y)=0$ and $b=1$ elsewhere. The only blank vertex is $y$, and $y$ is not
adjacent to $r$ in $B_0$. Write $N(y)\cap T\subseteq\{p,p'\}$. If $N(y)\cap T$ is empty, let $p$ be
any vertex of $T$. Perform:

\begin{center}\small
\begin{tabular}{llll}
\toprule
step & transfer & (supply, bound on $\bl$) & reason for the bound\\
\midrule
1 & $p\tr r$ (single) & $(1,0)$ & no blank neighbour of $r$ ($ry$ held)\\
2 & $(T-p)\tr r$ (grouped), donor $z_2$ & $(4,1)$ & blank neighbours of $r$: $p$\\
3 & $(T-\{p,z_2\})\tr r$, donor $z_3$ & $(3,2)$ & $p,z_2$\\
4 & $r\tr s_1$ & $(5,4)$ & $Z=\{p,z_2,z_3\}$, and possibly $y$\\
5 & $r\tr s_2$ & $(4,3)$ & $Z$ only\\
6 & $r\tr y$ & $(3,2)$ & $N(y)\cap Z\subseteq\{p,p'\}$\\
\bottomrule
\end{tabular}
\end{center}

Here $\{s_1,s_2\}=T-Z$, ordered so that $s_1=p'$ if $p'\in T-Z$. Since $p\in Z$, at most one of
$s_1,s_2$ is adjacent to the blank vertex $y$, and if one is, it is $s_1$. In steps 4--5 the blank
vertices are $Z$ and $y$, so $\bl(s_1)\le4$ and $\bl(s_2)\le3$. At step 6 the blank vertices are
$Z$ alone, and $y$ sees at most $p,p'$ of them. Every step satisfies Lemma~\ref{lem:fan}, all held
edges are restored, and the final counts are $b(r)=2$, $b(s_1)=b(s_2)=2$, $b=0$ on $Z$ and $b(y)=1$.

(ii) Hold $H=\{rt:t\in T-q\}\cup\{ra,rb,ab\}$. In $B_0$, $r$ has even degree $d(r)-6\ge2$, $q$ is
untouched and even (and not isolated), $T-q$, $a$ and $b$ are odd. Lemma~\ref{lem:seeds}(b) gives $b(r)=2$, $b(q)=0$,
and $b=1$ elsewhere. The blank vertex $q$ is adjacent neither to $r$ nor to $b$. Let $z_1$ be a
vertex of $N(a)\cap(T-q)$ if this set is nonempty, and any vertex of $T-q$ otherwise.

\begin{center}\small
\begin{tabular}{llll}
\toprule
step & transfer & (supply, bound) & reason\\
\midrule
1 & $b\tr r$ & $(1,0)$ & $q\notin N(r)$; $ra$ held\\
2 & $a\tr b$ & $(1,0)$ & $q\notin N(b)$; $a$ is now parked\\
3 & $z_1\tr r$ & $(1,0)$ & $a$, $q$ not adjacent to $r$\\
4 & $(T-q-z_1)\tr r$, donor $z_2$ & $(3,1)$ & $z_1$\\
5 & $r\tr s_1$ & $(5,4)$ & $Z=\{q,z_1,z_2\}$, and possibly $a$\\
6 & $r\tr s_2$ & $(4,3)$ & $Z$\\
7 & $r\tr a$ & $(3,2)$ & $|N(a)\cap Z|\le|N(a)\cap T|\le2$\\
\bottomrule
\end{tabular}
\end{center}

Here $\{s_1,s_2\}=T-Z$, and $N(a)\cap(T-q)\subseteq\{z_1,w\}$, so at most one receiver ($w$) is
adjacent to the parked vertex $a$. Order it first. The final counts are $b(r)=2$, two doubled
receivers and $Z$ blank.

(iii) Let $N_T(r)=\{t_1,t_2,t_3\}$ and hold $H=\{rt_1,rt_2,rt_3,ra,rb,ab\}$. In $B_0$, $r$ has odd
degree $d(r)-5$, $q_1,q_2$ are untouched and even (and not isolated), and the rest is odd. So
$\Ev(B_0)=\{q_1,q_2\}$, $B_0$ has no isolated vertex, and Lemma~\ref{lem:seeds}(b) gives $b(q_1)=2$, $b(q_2)=0$, $b(r)=1$ and $b=1$ elsewhere. Let $z_1$
be the vertex of $N(a)\cap N_T(r)$ if there is one, and any $t_i$ otherwise.

\begin{center}\small
\begin{tabular}{llll}
\toprule
step & transfer & (supply, bound) & reason\\
\midrule
1 & $b\tr r$ & $(1,0)$ & $q_1,q_2\notin N(r)$\\
2 & $a\tr b$ & $(1,0)$ & $q_2\notin N(b)$; $a$ parked\\
3 & $z_1\tr r$ & $(1,0)$ & \\
4 & $(N_T(r)-z_1)\tr r$, donor $z_2$ & $(2,1)$ & $z_1$\\
5 & $r\tr s$ & $(4,3)$ & $Z=\{q_2,z_1,z_2\}$; $s\notin N(a)$\\
6 & $r\tr a$ & $(3,2)$ & $N(a)\cap Z\subseteq\{q_2,z_1\}$\\
\bottomrule
\end{tabular}
\end{center}

Here $s$ is the third vertex of $N_T(r)$. It is not adjacent to $a$ because
$N(a)\cap N_T(r)\subseteq\{z_1\}$. The final counts are $b(r)=b(q_1)=b(s)=2$ and $b=0$ on $Z$.
\end{proof}

\begin{remark}
Part (i) overlaps six-even-incidence Thm.~1(1) when $|S(y)|\le2$, where $S(y)=N(y)\cap U$. That
theorem allows any odd vertex with one or two core neighbours, but its proof rests on the
unreplayed trees of Lemma~2. Part (i) is new for $|S(y)|=3$, i.e.\ $S(y)=\{r,p,p'\}$, for every
exterior degree. The four-exterior-neighbors analysis left the type $(|S|,e)=(3,4)$ open for four of
64 masks and $(3,\ge6)$ open in general. Parts (ii) and (iii) are, as far as we can tell, the first
closures specific to $D=K_2$ with $r$ on the edge and to the $P_3$, $K_3$ and $C_4$ root positions.
\end{remark}

\begin{theorem}[Two resources]\label{thm:tworesources}
Let $Q=\emptyset$ and $d(r)\ge10$. Suppose $r$ has distinct odd neighbours $y,a,b$ with $ab\in E$,
$y\notin N(a)\cup N(b)$, $|N(a)\cap T|\le4$ and $|N(y)\cap N(a)\cap T|\le2$, and $d(y)\ge3$. Then
$r$ passes. No bound on $|N(y)\cap T|$ is needed.
\end{theorem}
\begin{proof}
Hold $\{rt:t\in T\}\cup\{ry,ra,rb,ab\}$. Then $\Ev(B_0)=\{r,y\}$, with $r$ of degree
$d(r)-8\ge2$. The seed has $b(r)=2$ and $b(y)=0$. Let $z_1\in T-N(a)$, which exists since
$|N(a)\cap T|\le4$, and let $C=N(y)\cap N(a)\cap T$.

\begin{center}\small
\begin{tabular}{llll}
\toprule
step & transfer & (supply, bound) & reason\\
\midrule
1, 2 & $b\tr r$, $a\tr b$ & $(1,0)$, $(1,0)$ & $y\notin N(r)$ in $B_0$, $y\notin N(b)$\\
3 & $z_1\tr r$ & $(1,0)$ & \\
4 & $P_2\tr r$, donor $z_2$ & $(2\text{ or }4,1)$ & $P_2=C$ if $|C|=2$, else $P_2=T-z_1$\\
5 & $(T-\{z_1,z_2\})\tr r$ & $(3,2)$ & \\
6 & $r\tr s_1$ & $(6,5)$ & $Z$, plus possibly $y$ and $a$\\
7 & $r\tr s_2$ & $(5,4)$ & $Z$, plus at most one of $y,a$\\
8 & $r\tr y$ & $(4,3)$ & $|N(y)\cap Z|\le3$, and $y\notin N(a)$\\
9 & $r\tr a$ & $(3,2)$ & $z_1\in Z-N(a)$\\
\bottomrule
\end{tabular}
\end{center}

A receiver $s\in T-Z$ sees the parked $y$ and $a$ only if $s\in N(y)$, resp.\ $N(a)$. Both happen
only if $s\in C$. Since $z_1\notin N(a)\supseteq C$, we have $C\subseteq T-z_1$. If $|C|=2$,
step 4 offers $P_2=C$, so $|C\cap(T-Z)|\le1$. If $|C|\le1$, this holds trivially. Ordering the
receivers with the one in $C$ (if any) first gives the bounds in steps 6--7. Finally
$b(r)=2+1+3-2-2=2$.
\end{proof}

\begin{theorem}[Relay to a quiet vertex]\label{thm:relay}
Let $Q=\emptyset$ and $d(r)\ge8$. Suppose $r c_1c_2\cdots c_m a$ $(m\ge1)$ is a path of odd
vertices such that:
\begin{itemize}
\item $|N(c_1)\cap T|\le4$;
\item $N(c_i)\cap U=\emptyset$ for $2\le i\le m$;
\item $N(a)\cap U=\emptyset$, $a\notin N(r)$, and $a\notin N(c_i)$ for $i<m$;
\item $d(a)\ge3$.
\end{itemize}
Then $r$ passes.
\end{theorem}
\begin{proof}
Hold $\{rt:t\in T\}$ and the edges of the path. Then $\Ev(B_0)=\{r,a\}$, and the seed has
$b(r)=2$, $b(a)=0$. Choose $z_1\in T-N(c_1)$. Perform $z_1\tr r$ $(1,0)$, two grouped receipts at
$r$ $(4,1)$ and $(3,2)$, and $r\tr s_1$, $r\tr s_2$ $(5,3)$ and $(4,3)$. In these steps $a$ is the
only blank vertex outside $Z$. It is adjacent to no vertex of $U$, so the bounds hold. Then apply
the forward relay (Lemma~\ref{lem:relay}(a)) along $rc_1\cdots c_ma$. The first step is
$r\tr c_1$ with supply $3$ and $\bl(c_1)\le|N(c_1)\cap Z|\le2$, since $z_1\notin N(c_1)$ and $c_1a$
is held or $c_1\notin N(a)$. Every later receiver ($c_2,\dots,c_m$ and finally $a$) has no blank neighbour: these vertices have
no core neighbour, $a$ is adjacent to no $c_i$ with $i<m$, and $c_ma$ is held until the last step.
Finally $b(r)=2$.
\end{proof}

\begin{corollary}[Residual constraints at $D=\emptyset$]\label{cor:residual}
Let $r$ be a root that fails, in a graph with $B[U]=K_6$ (with or without full vertices), with
$d(r)\ge10$ and $\delta(B)\ge3$. Then:
\begin{enumerate}[label=(\alph*)]
\item every odd neighbour of $r$ has at least three neighbours in $T$, so $|S(y)|\ge4$ for
$y\in N(r)\cap\Od(B)$;
\item no odd neighbour $c$ of $r$ with $|N(c)\cap T|\le4$ is joined to a vertex with no core
neighbour by a path through vertices with no core neighbour, as in Theorem~\ref{thm:relay} (take
a shortest such path; then $a$ is adjacent to no $c_i$ with $i<m$, and $a\notin N(r)$ because $a$
has no core neighbour);
\item for every triangle $rab$ of odd vertices with $|N(a)\cap T|\le4$, every other odd neighbour
$y$ of $r$ is adjacent to $a$ or $b$ or satisfies $|N(y)\cap N(a)\cap T|\ge3$;
\item no vertex adjacent to all of $U$ satisfies the hypotheses of Theorem~\ref{thm:fullhub}.
\end{enumerate}
\end{corollary}

\section{Full hubs and batteries}\label{sec:fullhub}

\begin{definition}
Let $x\in\Od(B)$ be adjacent to all of $U$. A \emph{quiet battery} at $x$ is a triangle $xab$ with
$a,b\in\Od(B)$ and $N(a)\cap U=\emptyset$. The vertex $a$ is the \emph{parked} vertex. A
\emph{quiet pendant} at $x$ is an odd neighbour $y$ of $x$ with $N(y)\cap U=\emptyset$ and
$d(y)\ge3$.
\end{definition}

\begin{theorem}[Full hub]\label{thm:fullhub}
Let $B$ have even order, $U=\Ev(B)$, $|U|=6$, with any core $B[U]$, and let $x\in\Od(B)$ be
adjacent to all of $U$. Suppose $x$ carries quiet batteries $xa_1b_1$ and $xa_2b_2$ on five
distinct vertices with $a_1b_2\notin E$. Then every $r\in U$ passes. The same holds if $x$ carries
one quiet battery $xa_1b_1$ and a quiet pendant $y\notin\{a_1,b_1\}$ with $yb_1\notin E$, provided
$x$ has a neighbour outside $U\cup\{a_1,b_1,y\}$.
\end{theorem}
\begin{proof}
Two batteries. Hold $\{xu:u\in U\}$ and the six battery edges. In $B_0$ every vertex is odd:
$x$ lost ten edges, $a_i$ and $b_i$ lost two, and the vertices of $U$ lost one. The seed
(Lemma~\ref{lem:seeds}(a)) has $b=1$ everywhere. Fix $r\in U$.

\begin{center}\small
\begin{tabular}{llll}
\toprule
step & transfer & (supply, bound) & reason\\
\midrule
1, 2 & $b_1\tr x$, $a_1\tr b_1$ & $(1,0)$, $(1,0)$ & carry; $a_1$ parked\\
3, 4 & $b_2\tr x$, $a_2\tr b_2$ & $(1,0)$, $(1,0)$ & $xa_1$ held; $b_2\notin N(a_1)$; $a_2$ parked\\
5--7 & three grouped receipts $(U-r)\tr x$ & $(5,0)$, $(4,1)$, $(3,2)$ & spent donors; $a_i$ not adjacent to $x$\\
8--10 & $x\tr u$ for $u\in\{r\}\cup(T-Z)$ & $(6,3)$, $(5,3)$, $(4,3)$ & $|Z|=3$; the $a_i$ have no core neighbour\\
11 & $x\tr a_1$ & $(3,1)$ & $N(a_1)\cap Z=\emptyset$; possibly $a_2$\\
12 & $x\tr a_2$ & $(2,0)$ & \\
\bottomrule
\end{tabular}
\end{center}

The final counts are $b(x)=1+2+3-3-2=1$, $b=2$ on $r$ and on the two receivers in $T-Z$, $b=0$ on
$Z$, and $b=1$ on all odd vertices.

Pendant and battery. Hold $\{xu\}$, $xy$ and the battery. Now $x$ (degree $d(x)-9\ge2$) and $y$ are
the even vertices of $B_0$. Lemma~\ref{lem:seeds}(b) doubles $x$ and leaves $y$ blank. Carry on the
battery $(1,0)$, $(1,0)$, using $yb_1\notin E$. Then $b(x)=3$, and three receipts bring it to $6$. The donations and the paybacks to $y$ and $a_1$ have the same bounds as above, because
$y$ and $a_1$ have no core neighbour.
\end{proof}

\begin{remark}[Full vertices are not resources for the root hub]\label{rem:fullnot}
A full vertex used as the pendant $y$ in the proof of Theorem~\ref{thm:roothub}(i) has
$|N(y)\cap T|=5$. While parked it is a blank neighbour of every receiver in $T-Z$. This costs one
unit at each donation, which cancels the unit it supplies. So full vertices do not help the root
hub. This is why the residual of Corollary~\ref{cor:residual} consists of roots all of whose odd
neighbours are ``core-heavy''. The false-twin example $B=K_{10}$ minus the six edges from $9$ to $U$ satisfies
none of the hypotheses of \S\S\ref{sec:roothub}--\ref{sec:fullhub}. Every odd neighbour of $r$ is
full, and the only vertex with no core neighbour, $9$, supports a single battery. Yet $B$ passes,
by direct computation (every triple of $U$ is realisable).
\end{remark}

\section{All deletion degrees}\label{sec:general}

The counting of Remark~\ref{rem:capacity} is uniform in $d$. Let $|U|=d$ (even), $T=U-r$.

\begin{theorem}[Hubs at deletion degree $d$]\label{thm:general}
Let $B$ have even order and no isolated vertex, with $\Ev(B)=U$, $|U|=d\ge4$.
\begin{enumerate}[label=(\alph*)]
\item \emph{Root hub.} Let $r\in U$ be adjacent to all of $T$. Suppose $r$ has a quiet odd
neighbour $y$ ($N(y)\cap T=\emptyset$, $d(y)\ge3$). Suppose also that $r$ carries $\beta\ge d/2-3$
quiet battery triangles $ra_ib_i$ ($a_i,b_i$ odd, $N(a_i)\cap T=\emptyset$), all vertices
distinct, with $b_i\notin N(y)$ and $a_ib_j\notin E$ for $i<j$. Finally suppose $d(r)\ge d+2+2\beta$.
Then $r$ passes.
\item \emph{Full hub.} Let $x\in\Od(B)$ be adjacent to all of $U$ and carry $\beta\ge d/2-1$ quiet
batteries $xa_ib_i$ (distinct vertices, $N(a_i)\cap U=\emptyset$, $a_ib_j\notin E$ for $i<j$).
Then every $r\in U$ passes. The same holds with $\beta\ge d/2-2$ batteries and one quiet pendant
$y$, distinct from the battery vertices, with $b_i\notin N(y)$, if $x$ has a further neighbour.
\end{enumerate}
\end{theorem}

\begin{proof}
(a) Hold $\{rt:t\in T\}$, $ry$ and the battery triangles. Then $\Ev(B_0)=\{r,y\}$, since $r$ lost
$d+2\beta$ edges and the battery vertices lost two each. The seed has $b(r)=2$, $b(y)=0$. Carry
each battery: $b_i\tr r$, then $a_i\tr b_i$, with bounds $0$ because $y$ and the earlier parked
$a_j$ are not adjacent to $b_i$. Now $b(r)=2+\beta$. Receive $d/2$ times from $T$. The $j$-th
receipt ($0\le j<d/2$) has supply $d-1-j>j$ blank neighbours. Then $b(r)=2+\beta+d/2$ and
$|Z|=d/2$. Donate to the $d/2-1$ vertices of $T-Z$. The $i$-th donation has supply
$2+\beta+d/2-i>d/2\ge\bl$, because $i\le d/2-2\le\beta+1$ and the parked vertices have no
neighbour in $T$. Then $b(r)=\beta+3$. Pay back $y,a_1,\dots,a_\beta$. The $j$-th payback has
supply $\beta+3-j$ against at most $\beta-j$ parked blank neighbours. The final count is
$b(r)=2$.

(b) The same computation applies with seed $b(x)=1$ (Lemma~\ref{lem:seeds}(a)), $d/2$ receipts and
$d/2$ donations (to $r$ and to $T-Z$). The $i$-th donation has supply $1+\beta+d/2-i>d/2$ for
$i\le d/2-1\le\beta$. Then there are $\beta$ paybacks. With a pendant, the seed is $b(x)=2$ by
Lemma~\ref{lem:seeds}(b).
\end{proof}

\begin{remark}
For $d=4$ the root hub needs no battery and no quietness. The single donation has supply $4$
against at most $3$ blank neighbours, and the payback has supply $3$ against at most $2$. This
recovers T(4)-i at every universal root with at least three odd neighbours, one of which has
degree at least $3$ (so that it is not isolated in the seed graph), consistently with
Fan--Hou--Zhou. For $d=8$, (a) needs
one battery and (b) needs three batteries (or two and a pendant). These are, to our knowledge, the
first closures of the open test T(8)-i at roots adjacent to all other even vertices. The gap
between (a) and (b) is the gap of Remark~\ref{rem:capacity}: a root hub ends with two ends, a full
hub with one.
\end{remark}

\section{Full pairs do not fit the six-clique stages}\label{sec:stage}

Six-clique Thm.~1 and Cor.~5~\cite{6CL} use three stages. At each stage two clique vertices are blank, the
present clique edges form $D_0=\{23,45\}$, $D_1=D_0+\{02,13\}$ or
$D_2=D_1+\{14,01,03,35\}$, and Lemma~4 inserts all spokes of a pair of equal-type exterior vertices
while restoring the end vector. Cor.~5 requires $|A|<6$ so that some stage has a blank pair
meeting $A$ at most once.

\begin{proposition}\label{prop:stage}
At each of the three stages, with blanks $\{0,1\}$, $\{2,3\}$, $\{4,5\}$ and present clique edges
$D_0,D_1,D_2$ respectively, consider two full vertices $x,y$ with one end each and no spoke
present. No sequence of Fan transfers (Lemma~\ref{lem:fan}, grouped offers branching on every
donor, worst-case blank counting with exterior vertices never blank) inserts all twelve spokes
and restores all end counts. This holds whether or not $xy\in E$, and also when $x,y$ start with a
two-spoke seed row $xpy$ through any clique vertex $p$. The same search re-finds Lemma~4 for
five-contact pairs.
\end{proposition}
\begin{proof}
Exhaustive AND--OR search over all receivers and all inclusion-minimal donor sets; script
\texttt{scripts/stage\_search.py}, output \texttt{scripts/stage\_search.out}.
\end{proof}

So any extension of six-clique Thm.~1 or Thm.~6 to full vertices must interleave a full pair with
the core transfers across stages, or change the core sequence. Corollary~\ref{cor:nonlocal} says
the final step must then be non-local, as the final clique path of six-clique Thm.~1 is.

\section{What is closed and what remains}\label{sec:remains}

\subsection*{Closed} At every residual root position: the small-contact witnesses of
Theorem~\ref{thm:roothub}. At roots adjacent to all of $T$: also those of
Theorems~\ref{thm:tworesources} and~\ref{thm:relay}. For $D=\emptyset$ with full
vertices: any full vertex with two quiet batteries, or a quiet pendant and one quiet battery
(Theorem~\ref{thm:fullhub}). For every $d$: Theorem~\ref{thm:general}.

\subsection*{Not closed} Q2 itself. A failure at a root $r$ with $D=\emptyset$ and $d(r)\ge10$
must satisfy Corollary~\ref{cor:residual} as well as the earlier constraints: the bad-vertex list,
$\delta\ge6$, $3$-connectivity of $J$, and $|S(x)|\notin\{1,2\}$. In particular all odd neighbours
of $r$ are core-heavy ($|S|\ge4$), and vertices with no core neighbour are separated from $r$'s
non-$T$-full neighbours.

\subsection*{Where the method stops} Three limitations are visible.
\begin{enumerate}
\item \emph{Donor control.} Grouped receipts return an uncontrolled donor. With
$|N(y)\cap T|=3$ in Theorem~\ref{thm:roothub}(i), an adversarial choice of $z_2,z_3$ can always
defeat the payback. Full control of $Z$ would close $|N(y)\cap T|=3$.
\item \emph{No resources.} If every odd vertex within distance two of $r$ is core-heavy, the
counting in Remark~\ref{rem:capacity} leaves the root hub one unit short. The full hub is two
units short without quiet structure. The false-twin example is of this kind.
\item \emph{Inequality versus chain.} Fan's inequality is only sufficient. At equality Fan's
injection (from end occurrences at the reachable set into edges at the receiver) is a bijection,
and every blank neighbour of the receiver is reached. In the stuck state of the root hub
($b(r)=3$, receiver $s_2$, three blanks) this means the predecessor chains from the three
$r$-ending paths reach all of $Z$ through $s_2$, and every non-blank vertex they reach has exactly
one end. A structural use of
this rigidity, a genuine ``two-end transfer'', is the most promising missing ingredient.
\end{enumerate}

\begin{question}
In the apex class with $D=\emptyset$, does every root have, within distance two, an odd vertex
whose core neighbourhood is small enough to act as a pendant or a battery?
\end{question}

\subsection*{Observation (small computation)} In nine random graphs with $B[U]=K_6$ and one or two
full vertices (orders $10$ and $12$), \emph{every} 3-subset of $U$ was realisable, and so was the
false-twin example. This is consistent with uniformity of the endpoint matroid
(connections~\S M(a)) in this class. It is evidence only.

\section{Formalization}\label{sec:lean}
\sloppy

\begin{remark}[Formalization]\label{rem:lean}
Lemma~\ref{lem:relay}, Lemma~\ref{lem:seeds}, Theorem~\ref{thm:roothub},
Theorem~\ref{thm:fullhub} and Theorem~\ref{thm:general} are formalized in Lean~4 with Mathlib
(\texttt{lean/GallaiLean/Relay.lean}, \texttt{HubPhases.lean}, \texttt{D6Frontier.lean}), with no
\texttt{sorry}; \texttt{\#print axioms} shows only \texttt{propext}, \texttt{Classical.choice} and
\texttt{Quot.sound}. Input (A) is the Lean theorem \texttt{lovaszStatement}, so no result rests on
an unproved hypothesis. Every row of every step table is one application of \texttt{fan\_single}
or \texttt{fan\_grouped}, so the kernel checks each (supply, bound) pair.

\smallskip
\begin{center}\footnotesize
\begin{tabular}{ll}
\toprule
result & Lean name\\
\midrule
Lemma~\ref{lem:relay}(a), forward relay & \texttt{Ach.forwardRelay}\\
Lemma~\ref{lem:relay}(b), carry & \texttt{Ach.carry}\\
Lemma~\ref{lem:seeds}(a), (b) & \texttt{seed\_odd}, \texttt{seed\_pair}\\
Remark~\ref{rem:capacity}: grouped receipts, donations & \texttt{Ach.receiveMany}, \texttt{Ach.giveAll}\\
step-table rows (receipt, donation, payback, battery) & \texttt{hub\_take}, \texttt{hub\_receive}, \texttt{hub\_give},\\
 & \texttt{hub\_payback}, \texttt{hub\_pay\_one}, \texttt{hub\_batteries}\\
the common run of \S\S\ref{sec:fullhub}--\ref{sec:general} & \texttt{hub\_run}\\
Theorem~\ref{thm:roothub}(i) & \texttt{root\_hub\_six}, \texttt{root\_hub\_six\_passes}\\
Theorem~\ref{thm:roothub}(ii), (iii) & \texttt{root\_hub\_six\_Q1}, \texttt{root\_hub\_six\_Q2} (and \texttt{\_passes})\\
Theorem~\ref{thm:fullhub}, two batteries & \texttt{full\_hub\_six}\\
Theorem~\ref{thm:fullhub}, battery and pendant & \texttt{full\_hub\_six\_pendant}\\
Theorem~\ref{thm:general}(a) & \texttt{root\_hub\_uniform}, \texttt{root\_hub\_uniform\_passes}\\
Theorem~\ref{thm:general}(b), batteries & \texttt{full\_hub\_uniform}, \texttt{full\_hub\_uniform\_passes}\\
Theorem~\ref{thm:general}(b), pendant & \texttt{full\_hub\_pendant\_uniform},\\
 & \texttt{full\_hub\_pendant\_uniform\_passes}\\
\bottomrule
\end{tabular}
\end{center}
\smallskip

The statements are at the generality of the text. Theorem~\ref{thm:general} is stated for an
arbitrary $d=|U|$, with no lower bound (the proofs never use $d\ge4$, and $d$ is even by the
handshake lemma). $U$ is given by $v\in U\iff d(v)$ even, the order is even, and the conclusion
is the rooted test \texttt{PassesAt G r} (a partition into at most $|V|/2$ paths with $b(r)\ge2$).
Each theorem is first proved in the stronger form ``some $S\subseteq U$ with $r\in S$ and
$|S|=d/2$ is realisable'' (one end at every odd vertex, two on $S$, none on $U-S$), and
\texttt{passesAt\_of\_realisable} converts it. Batteries are a list of pairs $(a_i,b_i)$, carried
in list order, and ``$a_ib_j\notin E$ for $i<j$'' is a \texttt{List.Pairwise} hypothesis. The
carry is formalized under a weaker hypothesis than in Lemma~\ref{lem:relay}(b): a receiver may be
adjacent to a blank $v_0$, and only adjacency in $H$ matters.

Theorem~\ref{thm:roothub}(ii), (iii) are formalized with the hypothesis that the vertices of $Q$
are not isolated, which the seed needs (see the audit note
\texttt{research/audit/d6-frontier-lean-audit.md}).
Not formalized: Theorems~\ref{thm:terminal} and~\ref{thm:terminalW} (pure counting),
Theorems~\ref{thm:tworesources} and~\ref{thm:relay}
(each a finite composition of the lemmas above), Corollary~\ref{cor:residual}, and
Proposition~\ref{prop:stage} (an exhaustive search).
\end{remark}

\fussy
\section*{Computational checks}

All scripts are in \texttt{scripts/} and use \texttt{tools/pathnum2.py} only to compute seeds
(partitions with exact end vectors, which exist by Lemma~\ref{lem:seeds}). The certificates
themselves are executed by \texttt{scripts/fan\_engine.py}. This is an independent implementation
of Fan's chain construction, self-tested by \texttt{test\_engine.py} on $800$ random instances.
Every step asserts that the actual number of blank neighbours is at most the bound used in the
proof, and every output partition is verified edge by edge. In a grouped step a replay follows the
donor chosen by the engine; the argument for an arbitrary donor is the one in the proofs.
\texttt{scripts/run\_all.sh}
re-runs everything and regenerates the \texttt{.out} files.

\begin{center}\footnotesize
\begin{tabular}{lll}
\toprule
script & checks & result\\
\midrule
\texttt{terminal\_check.py} & Thms.~\ref{thm:terminal}, \ref{thm:terminalW}: all excess distributions & only hub profiles survive\\
\texttt{replay\_root\_hub.py} & Thm.~\ref{thm:roothub}(i), $|N(y)\cap T|\in\{0,1,2\}$, $n=12,14$ & 10/10 verified\\
\texttt{replay\_shapes.py} & Thm.~\ref{thm:roothub}(ii), (iii): $D=K_2,P_3,K_3,C_4$ & 8/8 verified\\
\texttt{replay\_two\_resources.py} & Thm.~\ref{thm:tworesources}, $|N(a)\cap T|=3$ & 4/4 verified\\
\texttt{replay\_relay.py} & Thm.~\ref{thm:relay}, $m=1,2$ & 5/5 verified\\
\texttt{replay\_full\_hub.py} & Thm.~\ref{thm:fullhub} (two batteries), every root & 3 graphs $\times$ 6 roots\\
\texttt{replay\_general\_d.py} & Thm.~\ref{thm:general}, $d=8$, root and full hub & 3/3 verified\\
\texttt{stage\_search.py} & Prop.~\ref{prop:stage} & no program; Lemma 4 re-found\\
\texttt{small\_observation.py} & \S\ref{sec:remains}, Observation & all triples realisable (10 graphs)\\
\bottomrule
\end{tabular}
\end{center}

\begin{thebibliography}{9}
\bibitem{ES} Endpoint sets of path decompositions at the Gallai budget, working draft
(\texttt{manuscripts/endpoint-structure}), Lemmas 1--4; formalized in
\texttt{lean/GallaiLean/Realisable.lean}.
\bibitem{Fan} G.~Fan, Path decompositions and Gallai's conjecture, J.\ Combin.\ Theory Ser.~B 93
(2005) 117--125, Lemmas 3.3--3.4; formalized in \texttt{lean/GallaiLean/Fan.lean}.
\bibitem{L68} L.~Lov\'asz, On covering of graphs, in: Theory of Graphs (Tihany 1966), Academic
Press, 1968, 231--236.
\bibitem{6CL} Prescribed path ends at a clique of six even vertices (six-clique-prescribed-ends),
November 2026; Lemmas 2--4, Thm.~1, Cor.~5, Thm.~6.
\bibitem{FTR} A false-twin return at a full exterior vertex, technical note, November 2026.
\end{thebibliography}

\end{document}
